\section{Analyse der Schwachstelle}
\label{sec:analyse}
Im Folgenden werden die Projektstruktur und die zugrundeliegende Sicherheitslücke analysiert.

\subsection{Projektstruktur}

Das Backend von dnchub-app ist in drei Schichten gegliedert.

\begin{itemize}
    \item Der API-Layer enthält die FastAPI-Router unter \texttt{app/api/v1/}.
    \item Der Service-Layer kapselt die Geschäftslogik in \texttt{app/shared/services/}.
    \item Der Model-Layer hält die SQLAlchemy-Modelle in \texttt{app/models/}.
\end{itemize}

Die Mandantentrennung hängt im gesamten Projekt an einer einzigen Spalte,
nämlich \texttt{organization\_id}. Sie wird in fast jeder
Geschäftsdaten-Tabelle geführt und soll bewirken, dass Anfragen immer nur
Zeilen der eigenen Organisation zurückliefern.

\subsection{Codeanalyse}

\subsubsection{Analyse der Klasse \texttt{BaseService}}

Die zentrale Klasse \texttt{BaseService} befindet sich in
\texttt{backend/app/shared/services/base.py} und implementiert sieben
Standard-CRUD-Operationen, die von allen Modulen des Systems geerbt werden.
Beim Durchsehen der Klasse fällt eine konsistente Inkonsistenz auf. Die
Listenmethoden (\texttt{get\_multi}, \texttt{count}) implementieren korrekt
sowohl einen Mandanten- als auch einen Soft-Delete-Filter, die
Einzelzugriffsmethoden (\texttt{get}, \texttt{get\_or\_404}, \texttt{delete}) hingegen nicht.

Tabelle \ref{tab:baseservice_audit} zeigt den Sicherheitsstatus aller Methoden.

\begin{table}[H]
    \begin{center}
        \renewcommand{\arraystretch}{1.4}
        \begin{tabular}{|l|c|c|}
            \hline
            \textbf{Methode}        & \textbf{IDOR-sicher}           & \textbf{Soft-Delete}           \\
            \hline
            \texttt{get()}          & \textcolor{red}{Nein}          & \textcolor{red}{Nein}          \\
            \texttt{get\_or\_404()} & \textcolor{red}{Nein}          & \textcolor{red}{Nein}          \\
            \texttt{get\_multi()}   & \textcolor{green!60!black}{Ja} & \textcolor{green!60!black}{Ja} \\
            \texttt{count()}        & \textcolor{green!60!black}{Ja} & \textcolor{green!60!black}{Ja} \\
            \texttt{create()}       &  --                            & --                             \\
            \texttt{update()}       &  --                            & --                             \\
            \texttt{delete()}       & \textcolor{red}{Nein}          & --                             \\
            \hline
        \end{tabular}
    \end{center}
    \caption{Sicherheitsaudit der Methoden in \texttt{BaseService}}
    \label{tab:baseservice_audit}
\end{table}

Listing \ref{lst:base_vulnerable} zeigt den verwundbaren Code. In \texttt{get()} wird
ausschliesslich nach \texttt{id} gefiltert, weder \texttt{organization\_id} noch
\texttt{deleted\_at} werden geprüft. Da \texttt{get\_or\_404()} intern \texttt{get()}
aufruft, erbt es beide Schwachstellen gleich mit.

\begin{lstlisting}[style=custompython,
  caption={Verwundbare Methoden in \texttt{base.py}},
  label={lst:base_vulnerable}]
def get(self, db: Session, id: str) -> ModelType | None:
    """Get a single record by ID."""
    # PROBLEM, kein organization_id-Filter, kein deleted_at-Filter
    result = db.execute(select(self.model).where(self.model.id == id))
    return result.scalar_one_or_none()

def get_or_404(self, db: Session, id: str) -> ModelType:
    """Get a single record by ID or raise NotFoundError."""
    # PROBLEM, delegiert an get() und erbt beide Schwachstellen
    obj = self.get(db, id)
    if obj is None:
        raise NotFoundError(self.model.__name__, id)
    return obj
\end{lstlisting}

Zum Vergleich zeigt Listing \ref{lst:base_correct} das Referenzmuster aus
\texttt{get\_multi()}, das beide Filter korrekt per \texttt{hasattr}-Guard implementiert.

\begin{lstlisting}[style=custompython,
  caption={Korrektes Referenzmuster aus \texttt{get\_multi()}},
  label={lst:base_correct}]
def get_multi(self, db, *, organization_id=None,
              include_deleted=False):
    query = select(self.model)
    # Korrektes Mandanten-Scoping mit Guard
    if organization_id and hasattr(self.model, "organization_id"):
        query = query.where(
            self.model.organization_id == organization_id)
    # Korrekter Soft-Delete-Filter mit Guard
    if not include_deleted and hasattr(self.model, "deleted_at"):
        query = query.where(self.model.deleted_at.is_(None))
    ...
\end{lstlisting}

Ein weiterer sicherheitsrelevanter Befund betrifft \texttt{delete()}, denn auch
diese Methode ruft intern das ungesicherte \texttt{get()} auf. Ein Angreifer
könnte damit nicht nur Datensätze fremder Mandanten lesen, sondern diese auch
soft-löschen, und das ohne jede \texttt{organization\_id}-Prüfung.

\subsubsection{Identifikation aller Aufrufer von \texttt{get\_or\_404()}}

Eine Volltextsuche im gesamten Backend (\texttt{grep -r "get\_or\_404"}) ergibt
83 Aufrufstellen in 29 Quelldateien. Tabelle \ref{tab:callers_summary}
fasst die Verteilung modulweise zusammen, wobei die Zeilen
\texttt{fuel*.py}, \texttt{tool\_*.py} und \texttt{sat/api/*.py} jeweils
mehrere gleichartige Dateien bündeln.

\begin{table}[H]
    \begin{center}
        \renewcommand{\arraystretch}{1.4}
        \begin{tabular}{|l|r|l|}
            \hline
            \textbf{Datei / Modul}                               & \textbf{Aufrufe} & \textbf{Handlungsbedarf} \\
            \hline
            \texttt{shared/api/organizations.py}                 & 5                & keiner                   \\
            \texttt{shared/api/users.py}                         & 2                & Pflicht-Fix              \\
            \texttt{shared/api/notifications.py}                 & 2                & Pflicht-Fix              \\
            \texttt{shared/api/documents.py}                     & 4                & Pflicht-Fix              \\
            \texttt{shared/api/employees.py}                     & 3                & Pflicht-Fix              \\
            \texttt{shared/api/cost\_centers.py}                 & 5                & Pflicht-Fix              \\
            \hline
            \texttt{modules/fleet/api/vehicles.py}               & 5                & Pflicht-Fix              \\
            \texttt{modules/fleet/api/maintenance.py}            & 7                & Pflicht-Fix              \\
            \texttt{modules/fleet/api/gps.py}                    & 7                & Pflicht-Fix              \\
            \texttt{modules/fleet/api/tickets.py}                & 4                & Pflicht-Fix              \\
            \texttt{modules/fleet/api/fuel.py}                   & 2                & Pflicht-Fix              \\
            \texttt{modules/fleet/api/fuel\_pumps.py}            & 2                & Pflicht-Fix              \\
            \texttt{modules/fleet/api/fuel\_pump\_deliveries.py} & 2                & Pflicht-Fix              \\
            \texttt{modules/fleet/services/fuel\_pump.py}        & 4                & Pflicht-Fix              \\
            \texttt{modules/fleet/services/fuel*.py}             & 2                & Pflicht-Fix              \\
            \hline
            \texttt{modules/tools/api/tools.py}                  & 3                & Pflicht-Fix              \\
            \texttt{modules/tools/api/tool\_cases.py}            & 4                & Pflicht-Fix              \\
            \texttt{modules/tools/api/consumables.py}            & 3                & Pflicht-Fix              \\
            \texttt{modules/tools/api/tool\_*.py}                & 7                & Pflicht-Fix              \\
            \texttt{modules/tools/services/tool\_assignment.py}  & 1                & Pflicht-Fix              \\
            \hline
            \texttt{modules/sat/api/*.py}                        & 9                & Pflicht-Fix              \\
            \hline
            \textbf{Total}                                       & \textbf{83}      &                          \\
            \hline
        \end{tabular}
    \end{center}
    \caption{Verteilung der \texttt{get\_or\_404()}-Aufrufstellen nach Modul.}
    \label{tab:callers_summary}
\end{table}

Aus der Analyse ergeben sich zwei Stufen von Handlungsbedarf, die in
der letzten Spalte der Tabelle bereits eingetragen sind.

\begin{itemize}
    \item \textbf{Keiner (5 Aufrufstellen).} \texttt{organizations.py}
          greift auf das \texttt{Organization}-Modell zu, das das
          Wurzelelement der Mandantenhierarchie ist und selbst kein
          \texttt{organization\_id}-Attribut besitzt. Endpunkte sind
          entweder Admin-only oder lesen die eigene Organisation via
          \texttt{current\_user.\allowbreak organization\_id}, sodass
          kein IDOR-Risiko besteht.

    \item \textbf{Pflicht-Fix (78 Aufrufstellen).} Alle übrigen
          Aufrufer verarbeiten mandantenspezifische Ressourcen und
          übergeben dem Aufruf kein \texttt{organization\_id}-Argument.
          Jeder dieser Aufrufe muss auf die neue Signatur umgerüstet
          werden und
          \texttt{current\_user.\allowbreak organization\_id}
          durchreichen. Die beiden Admin-Endpunkte in
          \texttt{users.py} fallen ebenfalls in diese Gruppe, weil ein
          Admin der Organisation~A sonst Benutzer anderer Mandanten
          lesen, ändern oder löschen kann, siehe Abschnitt
          \ref{sec:befunde-ausserhalb-baseservice}.
\end{itemize}

Auffällig ist zusätzlich, dass in den Service-Klassen
\texttt{fuel\_pump.py} und \texttt{tool\_assignment.py}
\texttt{self.get\_or\_404()} intern genutzt wird. Der Aufruf erfolgt
also nicht aus dem API-Layer, sondern aus einem anderen Service,
weshalb \texttt{organization\_id} durch den Service-Aufrufer
durchgereicht werden muss.


\subsubsection{Modelle mit \texttt{organization\_id} und \texttt{deleted\_at}}

Für die korrekte Implementierung des Fixes muss bekannt sein, welche
Modelle welche Felder tragen, da \texttt{BaseService} generisch über alle
Modelltypen operiert (vgl. die ORM-Abfrage-Semantik in
\parencite{sqlalchemy_docs}). Konkret wurde mit
\texttt{grep -rE "organization\_id|deleted\_at" backend/app/models/} jede
SQLAlchemy-Klasse darauf geprüft, ob sie \texttt{organization\_id},
\texttt{deleted\_at} oder beide Felder deklariert. Das Ergebnis wurde manuell
mit dem Klassennamen kategorisiert. Die Analyse
aller 31 Modell-Klassen ergibt die Klassifikation in Tabelle
\ref{tab:model_matrix}.

\begin{table}[H]
    \begin{center}
        \renewcommand{\arraystretch}{1.3}
        \begin{tabular}{|l|c|c|l|}
            \hline
            \textbf{Modell-Klasse}        & \textbf{org\_id}               & \textbf{deleted\_at}           & \textbf{Anzuwendende Filter} \\
            \hline
            \multicolumn{4}{|l|}{\textit{Modul fleet}}                                                                                     \\
            \hline
            Vehicle, VehicleGroup         & \textcolor{green!60!black}{Ja} & \textcolor{green!60!black}{Ja} & beide                        \\
            FuelEntry, FuelPump           & \textcolor{green!60!black}{Ja} & \textcolor{green!60!black}{Ja} & beide                        \\
            FuelPumpDelivery              & \textcolor{green!60!black}{Ja} & \textcolor{green!60!black}{Ja} & beide                        \\
            MaintenanceTask / Schedule    & \textcolor{green!60!black}{Ja} & \textcolor{green!60!black}{Ja} & beide                        \\
            Ticket, Geofence              & \textcolor{green!60!black}{Ja} & \textcolor{green!60!black}{Ja} & beide                        \\
            Trip                          & \textcolor{green!60!black}{Ja} & \textcolor{red}{Nein}          & nur org\_id                  \\
            GpsAlert                      & \textcolor{green!60!black}{Ja} & \textcolor{red}{Nein}          & nur org\_id                  \\
            VehicleGroupMember            & \textcolor{red}{Nein}          & \textcolor{red}{Nein}          & keine                        \\
            TripPosition                  & \textcolor{red}{Nein}          & \textcolor{red}{Nein}          & keine                        \\
            \hline
            \multicolumn{4}{|l|}{\textit{Modul tools}}                                                                                     \\
            \hline
            Tool, ToolCase, ToolCategory  & \textcolor{green!60!black}{Ja} & \textcolor{green!60!black}{Ja} & beide                        \\
            ToolLocation, Consumable      & \textcolor{green!60!black}{Ja} & \textcolor{green!60!black}{Ja} & beide                        \\
            ToolAssignment                & \textcolor{green!60!black}{Ja} & \textcolor{red}{Nein}          & nur org\_id                  \\
            ToolCalibration               & \textcolor{green!60!black}{Ja} & \textcolor{red}{Nein}          & nur org\_id                  \\
            \hline
            \multicolumn{4}{|l|}{\textit{Modul sat}}                                                                                       \\
            \hline
            SatAssistance, SatMachine     & \textcolor{green!60!black}{Ja} & \textcolor{green!60!black}{Ja} & beide                        \\
            SatCustomer, SatContact       & \textcolor{green!60!black}{Ja} & \textcolor{green!60!black}{Ja} & beide                        \\
            SatInterventionReport         & \textcolor{green!60!black}{Ja} & \textcolor{green!60!black}{Ja} & beide                        \\
            SatAttachment, SatServiceType & \textcolor{green!60!black}{Ja} & \textcolor{green!60!black}{Ja} & beide                        \\
            SatSpecialization             & \textcolor{green!60!black}{Ja} & \textcolor{green!60!black}{Ja} & beide                        \\
            SatEmployeeSpecialization     & \textcolor{red}{Nein}          & \textcolor{red}{Nein}          & keine                        \\
            \hline
            \multicolumn{4}{|l|}{\textit{Modul shared}}                                                                                    \\
            \hline
            Employee, User                & \textcolor{green!60!black}{Ja} & \textcolor{green!60!black}{Ja} & beide                        \\
            Document, CostCenter          & \textcolor{green!60!black}{Ja} & \textcolor{green!60!black}{Ja} & beide                        \\
            CostAllocation                & \textcolor{green!60!black}{Ja} & \textcolor{red}{Nein}          & nur org\_id                  \\
            Organization                  & \textcolor{red}{Nein}          & \textcolor{green!60!black}{Ja} & keine*                       \\
            Notification                  & \textcolor{red}{Nein}          & \textcolor{red}{Nein}          & keine                        \\
            NotificationPreferences       & \textcolor{red}{Nein}          & \textcolor{red}{Nein}          & keine                        \\
            \hline
        \end{tabular}
    \end{center}
    \caption{Modell-Klassifikation nach Feldern.}
    \label{tab:model_matrix}
\end{table}

Aus der Tabelle ergeben sich folgende Schlussfolgerungen für die
Implementierung.

\begin{itemize}
    \item \textbf{Beide Filter aktiv (27 Modelle).} Sowohl
          \texttt{organization\_id} als auch \texttt{deleted\_at} müssen in
          \texttt{get()} und \texttt{get\_or\_404()} berücksichtigt werden.
    \item \textbf{Nur \texttt{organization\_id} aktiv (5 Modelle).} Diese
          Modelle besitzen kein \texttt{deleted\_at}-Feld, weshalb Issue~\#14
          für sie nicht zutrifft. Die IDOR-Lücke aus Issue~\#5 bleibt aber
          relevant und wird über den Mandantenfilter geschlossen.
    \item \textbf{Keine Filter (5 Klassen).} Join-Tabellen wie
          \texttt{VehicleGroupMember}, benutzergebundene Tabellen wie
          \texttt{Notification} und die Root-Entity \texttt{Organization}
          tragen keinen Mandantenschlüssel und sind zudem nicht über das
          generische \texttt{get\_or\_404()} exponiert.
\end{itemize}

Das \texttt{hasattr}-Guard-Muster aus \texttt{get\_multi()} und \texttt{count()}
behandelt alle Fälle korrekt, weil es prüft, ob das jeweilige Modell das
gesuchte Attribut überhaupt besitzt. Dasselbe Muster, kann direkt in \texttt{get()} und \texttt{get\_or\_404()} übernommen werden.


\subsubsection{Befunde ausserhalb von \texttt{BaseService}}
\label{sec:befunde-ausserhalb-baseservice}

Der Wurzel-Fix greift nur dort, wo die Basisklasse tatsächlich
durchlaufen wird. Deshalb wurden alle Router-Dateien gezielt darauf
durchgesehen, ob sie \texttt{BaseService} umgehen oder nach dem
\texttt{get\_or\_404}-Lookup den nötigen Eigentümercheck weglassen.
Dabei fanden sich sechs eigenständige Lückenklassen.

\paragraph{Notifications ohne Eigentümer-Check.}
Notifications sind nicht an eine Organisation, sondern über
\texttt{user\_id} an einen einzelnen Benutzer gebunden. In GET, DELETE
und der mark-as-read-Route fehlt der Vergleich
\texttt{notification.user\_id == current\_user.id}, sodass jeder
authentifizierte Benutzer fremde Notifications lesen, als gelesen
markieren und löschen kann.

\paragraph{SAT-Assistances umgehen \texttt{BaseService}.}
\texttt{Assistance\allowbreak Service.get\_detail()} baut eigene
SQLAlchemy-Filter, ohne \texttt{BaseService.get()} aufzurufen. Dort
fehlt der Mandantenfilter vollständig, ein Angreifer kann also
Assistances fremder Organisationen abfragen.

\paragraph{Admin-Endpunkte in \texttt{users.py}.}
Die Routen \texttt{GET /users/\{user\_id\}},
\texttt{PATCH /users/\{user\_id\}} und
\texttt{DELETE /users/\{user\_id\}} rufen \texttt{get\_or\_404} und
\texttt{delete()} ohne \texttt{organization\_id} auf. Ein Admin aus
Organisation~A kann damit Benutzer anderer Mandanten lesen, ändern und
sogar löschen, was gravierender ist als die übrigen Befunde, weil
gleich drei HTTP-Verben betroffen sind.

\paragraph{DELETE in \texttt{sat/api/attachments.py}.}
\texttt{SatAttachmentService} erbt von \texttt{BaseService} und nimmt
in \texttt{delete()} bereits einen \texttt{organization\_id}-Parameter
entgegen. Der Router übergibt ihn jedoch nicht, weshalb jeder
SAT-Manager fremde Anhänge per UUID löschen kann. Der Patch ist klein,
die Auswirkung aber direkt Cross-Tenant.

\paragraph{GPS-Koordinate in fremde Trips.}
\texttt{POST /trips/\{trip\_id\}/\allowbreak position} schreibt eine
Koordinate in einen beliebigen Trip, ohne zu prüfen, ob dieser dem
eigenen Mandanten gehört. Fremde Trips lassen sich auf diesem Weg mit
falschen Positionsdaten verunreinigen. Gelesen wird dabei nichts, der
Schaden liegt in der Datenintegrität der fremden Flotte.

\paragraph{FK-List-Methoden ohne Mandantenfilter.}
Die zahlenmässig grösste Klasse umfasst 19 Service-Methoden quer durch
die Module \texttt{fleet}, \texttt{tools}, \texttt{sat} und
\texttt{shared}. Sie filtern List-Abfragen ausschliesslich nach einer
Fremdschlüssel-ID, ohne \texttt{organization\_id} in die
\texttt{WHERE}-Klausel aufzunehmen. Wer einen einzelnen
Fremdschlüsselwert mitliest oder errät, erhält damit die vollständige
Liste eines fremden Mandanten zurück. Weil der Angriff einen gültigen
Fremdschlüssel voraussetzt, ist diese Klasse weniger akut als die
direkten Cross-Tenant-Zugriffe weiter oben, aber wegen ihrer Breite
trotzdem relevant.


\subsubsection{Analyse der bestehenden Tests}

Die vollständige Test-Suite des Projekts liegt unter \texttt{backend/tests/} und ist
in drei Schichten gegliedert.

\begin{itemize}
    \item Die Unit- und API-Tests (\texttt{tests/test\_*.py}) prüfen
          einzelne Endpunkte mit echten Datenbank-Transaktionen.
    \item Die Integrationstests (\texttt{tests/integration/}) prüfen
          vollständige Use-Cases über mehrere CRUD Endpunkte hinweg.
    \item Die Smoke-Tests (\texttt{tests/smoke/}) führen
          Status-Code-Checks gegen eine Live-ähnliche DB aus. Bekannte
          Bugs sind mit \texttt{@pytest.mark.xfail} markiert.
\end{itemize}

\paragraph{Test-Fixtures (conftest.py).}
Alle Tests verwenden eine gemeinsame Fixture-Hierarchie. \texttt{db\_session}
rollt jede Transaktion via Savepoint zurück, \texttt{test\_organization} legt
eine Testorganisation an und \texttt{test\_admin\_user} sowie
\texttt{test\_fleet\_manager} sind User dieser Organisation. Die
Header-Fixtures \texttt{admin\_headers} und \texttt{manager\_headers}
erzeugen JWT-Bearer-Tokens via
\texttt{create\_access\_token()}. Ein
zweites Mandanten-Objekt existiert in den bestehenden Testfixtures
nicht. Genau dort setzen die nachfolgenden Proof-of-Concept-Versuche
(PoC) an.

\paragraph{Befund in integration/test\_vehicle\_crud.py.}
Die Datei enthält den Test
\texttt{test\_delete\_then\_get\_still\_returns\_soft\_deleted}, der nach
einem Soft-Delete explizit \texttt{status\_code == 200} erwartet und mit
dem Kommentar ``GET by ID still returns 200 (soft-deleted records are not
filtered in get\_or\_404)'' versehen ist, ein Eingeständnis des bekannten
Bugs. Dieser Test muss nach dem Patch auf \texttt{status\_code == 404}
umgestellt werden.

\paragraph{Befund in smoke/test\_tools.py.}
\texttt{test\_delete\_tool\_nonexistent\_returns\_204} bestätigt ein weiteres Symptom.
\texttt{DELETE} auf eine nicht existierende ID liefert \texttt{204} statt \texttt{404},
weil \texttt{delete()} intern \texttt{get()} ohne Soft-Delete-Filter aufruft.

\paragraph{Fehlender Cross-Tenant-Test.}
Keiner der bestehenden Tests prüft, ob ein User die Ressourcen einer anderen
Organisation lesen kann. Alle Fixtures operieren innerhalb derselben
\texttt{test\_organization}. Diese Lücke wird im nachfolgenden Abschnitt durch
dedizierte PoC-Tests geschlossen.

\subsection{Proof-of-Concept}
\label{sec:poc}

Bevor am Code etwas verändert wird, muss die Schwachstelle
reproduzierbar gezeigt werden. Der folgende Ablauf verwendet
\texttt{curl} gegen das lokal laufende Backend und lässt sich mit
wenigen Kommandozeilen-Aufrufen nachstellen. Das Backend wird unter
\texttt{http://localhost:8000} erwartet, \texttt{python3} dient nur
dazu, die JSON-Antworten lesbar zu machen. Die ausführliche Variante
mit Screenshots findet sich in Anhang~\ref{app:reproduktion}.

\subsubsection{Zwei unabhängige Mandanten registrieren.}
Die öffentliche \texttt{/auth/register}-Route legt neue Organisationen
samt erstem Benutzer an. Sie liefert einen JWT als
\texttt{access\_token} zurück.

\begin{lstlisting}[style=custombash, caption={Registrierung des Opfer- und des Angreifer-Mandanten.}]
# Opfer registrieren und Token in TOKEN_A ablegen
TOKEN_A=$(curl -s -X POST http://localhost:8000/api/v1/auth/register \
  -H "Content-Type: application/json" \
  -d '{"organization_name":"Opfer GmbH","email":"opfer@test.ch",
       "password":"opfer1234","first_name":"Opfer","last_name":"Opfer"}' \
  | python3 -c 'import json,sys; print(json.load(sys.stdin)["access_token"])')

# Angreifer registrieren und Token in TOKEN_B ablegen
TOKEN_B=$(curl -s -X POST http://localhost:8000/api/v1/auth/register \
  -H "Content-Type: application/json" \
  -d '{"organization_name":"Angreifer AG","email":"angreifer@test.ch",
       "password":"evil1234","first_name":"Angreifer","last_name":"Angreifer"}' \
  | python3 -c 'import json,sys; print(json.load(sys.stdin)["access_token"])')
\end{lstlisting}

\subsubsection{Ressource anlegen.}
Mit dem Token des Opfers wird ein Fahrzeug erzeugt. Die UUID des neuen
Datensatzes wird für die späteren Angriffsaufrufe benötigt und landet in
der Variable \texttt{VID}.

\begin{lstlisting}[style=custombash, caption={Das Opfer erzeugt ein Fahrzeug im eigenen Mandanten.}]
VID=$(curl -s -X POST http://localhost:8000/api/v1/vehicles \
  -H "Authorization: Bearer $TOKEN_A" \
  -H "Content-Type: application/json" \
  -d '{"registration_plate":"OPFER-001","make":"Toyota","model":"Camry",
       "year":1992,"type":"suv","fuel_type":"gasoline"}' \
  | python3 -c 'import json,sys; print(json.load(sys.stdin)["id"])')
\end{lstlisting}

\subsubsection{Cross-Tenant-GET (Issue~\#5).}
Der Angreifer ruft mit seinem eigenen JWT den Fahrzeug-Endpunkt unter
der UUID des Opfers auf. Korrektes Verhalten wäre ein \texttt{404 Not
Found}, weil das Fahrzeug seinem Mandanten nicht gehört. Vor dem Patch
antwortet der Server jedoch mit \texttt{200 OK} und liefert den
vollständigen Datensatz inklusive \texttt{organization\_id} des Opfers
aus. Damit ist die IDOR-Schwachstelle belegt.

\begin{lstlisting}[style=custombash, caption={IDOR per API. Vor dem Patch \texttt{200 OK} mit Opferdaten, nach dem Patch \texttt{404 Not Found}.}]
curl -s -X GET http://localhost:8000/api/v1/vehicles/$VID \
  -H "Authorization: Bearer $TOKEN_B" \
  -w "\nHTTP-Status: %{http_code}\n" \
  | python3 -m json.tool
\end{lstlisting}

\subsubsection{Soft-Delete-Bypass (Issue~\#14).}
Das Opfer löscht sein Fahrzeug regulär per \texttt{DELETE}. Der Server
setzt \texttt{deleted\_at}, gibt \texttt{204 No Content} zurück und
blendet den Datensatz aus der Liste aus. Die direkte Einzelabfrage per
UUID liefert den Datensatz allerdings weiterhin, inklusive gesetztem
\texttt{deleted\_at}-Feld. Korrektes Verhalten wäre auch hier ein
\texttt{404}. Damit ist der Soft-Delete-Bypass reproduziert.

\begin{lstlisting}[style=custombash, caption={Soft-Delete-Bypass: Der Datensatz wird gelöscht, aber per GET weiterhin ausgeliefert.}]
# Opfer löscht sein eigenes Fahrzeug regulär
curl -s -X DELETE http://localhost:8000/api/v1/vehicles/$VID \
  -H "Authorization: Bearer $TOKEN_A" \
  -w "\nHTTP-Status: %{http_code}\n"

# Vor dem Patch: 200 OK mit gesetztem deleted_at-Feld (Bypass)
# Nach dem Patch: 404 Not Found
curl -s -X GET http://localhost:8000/api/v1/vehicles/$VID \
  -H "Authorization: Bearer $TOKEN_A" \
  -w "\nHTTP-Status: %{http_code}\n" \
  | python3 -m json.tool
\end{lstlisting}

\subsubsection{Auswertung.}
Beide Angriffe lassen sich in weniger als einer Minute Terminalzeit
nachstellen und produzieren reproduzierbar die erwarteten
Fehl-Antworten. Damit ist der Angriffspfad konkret belegt und es kann
zur Behebung am Code übergegangen werden. Die gleichen Schritte werden
in Kapitel~\ref{sec:testing} als automatisierte Red/Green-Testsuite
gefasst, sodass der Patch nicht nur behoben, sondern auch gegen
Regressionen abgesichert ist.
